\documentclass[11pt]{article}

% Reusable, minimal assignment layout.
\usepackage{assignment}

% Your shared math macros.
\input{math_commands.tex}

% ------------------------------------------------------------------
% Edit these four lines for each assignment.
% Leave \course{} and/or \term{} empty if you do not want them shown.
% ------------------------------------------------------------------
\course{}
\term{}
\assignment{Problem Set 1}
\duedate{September 27, 2026, at 11:59 p.m.}

\begin{document}
\makeassignmenttitle

\begin{problem}[20]
Consider a venture that succeeds with probability $1/2$. For each dollar invested,
you earn a profit of \$1.60 if the venture succeeds and lose the invested dollar if it fails.
You may invest repeatedly, and each time you invest exactly half of your current wealth.
Starting with \$1,000,000:

\begin{subproblems}
  \item \points{5} What is the effect of combining this investment opportunity with the strategy
  of always investing half of your current wealth?

  \item \points{5} Do you gain or lose money on average?

  \item \points{10} If you make 10,000 independent investments of this kind, what is your
  expected wealth after the final investment? What is the probability that your final
  wealth is less than the \$1,000,000 with which you started?
\end{subproblems}
\end{problem}

\begin{problem}[10]
A standard 52-card deck is well shuffled, and five cards are dealt without replacement.
Find the probability of each of the following mutually exclusive poker-hand categories:
royal flush; other straight flush; four of a kind; full house; flush (excluding straight flushes);
straight (excluding straight flushes); three of a kind; two pairs; one pair; and high card only
(none of the preceding categories). Show how you count the number of hands in each category,
and verify that your counts sum to $\binom{52}{5}$.
\end{problem}

\begin{problem}[10]
Let $\mathcal{F}$ be a $\sigma$-field on $\Omega$. Prove that $\mathcal{F}$ is closed under
intersections and set differences; that is, for any $A,B\in\mathcal{F}$, show that
\[
  A\cap B\in\mathcal{F}
  \qquad\text{and}\qquad
  A\setminus B\in\mathcal{F}.
\]
\end{problem}

\begin{problem}[10]
Suppose that we are given a measure space $(\Omega,\mathcal{F},\mu)$. Show that whenever
$A\subseteq B$, we have
\[
  \mu(A)\leq \mu(B).
\]
Make explicit at each step where the defining properties of a measure and a $\sigma$-field are used.
\end{problem}

\begin{problem}[20]
Let $\Omega=\{R,M,L\}$ and take the $\sigma$-field $\mathcal{F}=2^\Omega$. Let $P$ be a
probability measure on $(\Omega,\mathcal{F})$ such that
\[
  P(\{R\})=\frac12,
  \qquad
  P(\{M\})=\frac14,
  \qquad
  P(\{L\})=\frac14.
\]
Define random variables $X$ and $Y$ by, for every $\omega\in\Omega$,
\[
  X(\omega)=2\,\mathbf{1}\{\omega=R\}+3\,\mathbf{1}\{\omega\in\{M,L\}\},
\]
\[
  Y(\omega)=2\,\mathbf{1}\{\omega\in\{M,L\}\}+3\,\mathbf{1}\{\omega=R\}.
\]

\begin{subproblems}
  \item \points{10} Find the cumulative distribution function of $X$ and sketch its graph.

  \item \points{10} Find $\sigma(X)$, the $\sigma$-field generated by $X$.
\end{subproblems}
\end{problem}

\begin{problem}[20]
Suppose that we are given a measurable space $(\Omega,\mathcal{F})$.

\begin{subproblems}
  \item \points{10} Suppose that $f:\Omega\to\mathbb{R}^d$ is measurable and
  $g:\mathbb{R}^d\to\mathbb{R}^r$ is Borel measurable; that is, for every
  $B\in\mathcal{B}(\mathbb{R}^r)$,
  \[
    g^{-1}(B)\in\mathcal{B}(\mathbb{R}^d).
  \]
  Show that $g\circ f$ is measurable.

  \item \points{10} Let $X:\Omega\to\mathbb{R}$ be a random variable and let
  $f:\mathbb{R}\to\mathbb{R}$ be Borel measurable. Prove that the $\sigma$-field generated by
  $X$ contains the $\sigma$-field generated by $f(X)$; that is, show that
  \[
    \sigma(f(X))\subseteq\sigma(X).
  \]
\end{subproblems}
\end{problem}

\begin{problem}[10]
Give an example of a probability space $(\Omega,\mathcal{F},P)$ and two random variables $X$
and $Y$ such that $X$ and $Y$ have the same distribution but are not the same random variable.
\end{problem}

\end{document}
